% ------------------------------------------------------------------------------
% Abstract
% ------------------------------------------------------------------------------
\section*{Abstract}
\addcontentsline{toc}{section}{Abstract}

\textbf{Background.} Adjuvant chemotherapy with levamisole and
5-fluorouracil (5-FU) after resection of stage B/C colon carcinoma was
established by the North Central Cancer Treatment Group (NCCTG) program in
the late 1980s, and the resulting individual patient dataset has become one
of the most widely used teaching resources in survival analysis. This
project asks what a modern biostatistician can extract from the same trial
today, when equipped with contemporary methods for missing data, sequential
monitoring, simulation, machine learning, Bayesian inference, and
heterogeneous treatment effects.

\smallskip
\noindent\textbf{Methods.} Among 929 patients randomized to observation
($n=315$), levamisole ($n=310$), or levamisole plus 5-FU ($n=304$), a
composite disease-free survival (DFS) endpoint (recurrence or death; 506
events) was reconstructed from the published event-type coding, and overall
survival (452 deaths) served as the secondary endpoint. The re-analysis
proceeds in six modules: (i) Kaplan--Meier estimation, log-rank tests,
Cox models with multiple imputation ($m=20$), competing-risk and
restricted-mean-survival-time sensitivity analyses, proportional-hazards
diagnostics, and bootstrap absolute
benefit; (ii) a redesign of the trial as a three-look group-sequential
design with O'Brien--Fleming efficacy spending and non-binding beta-spending
futility boundaries, including a re-enactment of interim monitoring on the
real event times; (iii) Monte Carlo simulation studies quantifying the
efficiency of covariate adjustment, the operating characteristics of the
sequential design, and confidence-interval coverage under non-proportional
hazards; (iv) cross-validated survival prediction with five
machine-learning models; (v) a Bayesian Weibull proportional-hazards
analysis with skeptical, weakly-informative, and enthusiastic priors,
sequential monitoring, and posterior predictive checks; and (vi) doubly
robust estimation of average and conditional treatment effects on the
five-year DFS event risk using jackknife pseudo-observations.

\smallskip
\noindent\textbf{Results.} Levamisole plus 5-FU reduced the hazard of the
composite DFS endpoint by 38\% (adjusted hazard ratio $0.621$, 95\% CI
$0.496$--$0.776$), with consistent evidence across unadjusted,
imputation-adjusted, and Fine--Gray competing-risk analyses, and an absolute
benefit of 16.8 percentage points at five years (59.2\% vs.\ 42.4\% DFS;
number needed to treat $=6$), and a gain of 0.63 disease-free years in
restricted mean survival time within five years. Under the redesigned sequential plan, the
first interim analysis after 118 events would have yielded a hazard ratio of
$0.52$ ($Z=3.38$), crossing the efficacy boundary of $2.96$; Bayesian
monitoring with a skeptical prior would have reported
$\Pr(\mathrm{HR}<1)=0.998$ after only 25\% of events. In simulation,
covariate adjustment increased power from 72.6\% to 80.9\% for a moderate
effect without inflating the type-I error, and directly illustrated the
non-collapsibility of the hazard ratio. Random survival forests achieved
the best validated discrimination (cross-validated $C=0.646$), and doubly
robust estimation located the largest absolute benefits in the
highest-risk patients.

\smallskip
\noindent\textbf{Conclusions.} A single 1980s trial, re-analyzed with
modern biostatistical machinery, supports a robust treatment effect and
illustrates how contemporary design and monitoring standards would have
accelerated its conclusion. The project demonstrates the complementary value
of frequentist, Bayesian, and machine-learning perspectives within one
coherent, fully scripted pipeline. All code, data, and outputs are
reproducible end-to-end from the accompanying repository.

\smallskip
\noindent\textbf{Keywords:} randomized clinical trial; survival analysis;
group-sequential design; multiple imputation; competing risks; Bayesian
inference; machine learning; heterogeneous treatment effects.
